% GENERATED FILE. Edit canonical lessons, then run npm run generate:books.
% canonical-source-hash: fnv1a64:f87d1c57a285f346
% canonical-lessons: HI-C07-haan, HI-C07-nahin

\chapter{Yes and No}
\label{ch:hi-yes-and-no}

\begin{chapteropening}
\textbf{By the end of this chapter:} \emph{I can say yes and no in Hindi and use \hi{नहीं} to negate a whole sentence.}

Say \hi{नहीं}, put it in front of a verb — \hi{मैं} \hi{नहीं} \hi{जानता} — and hear the single PIE *ne behind no, non and nein.
\end{chapteropening}

\section[\hi{हाँ}]{\hi{हाँ} (hāṃ) — yes}

\label{lesson:HI-C07-haan}



\begin{quote}
A single open vowel, said through the nose — that's Hindi's yes.
\end{quote}

\subsection*{You'll want to know — The word}
\textbf{\hi{हाँ}} = \textbf{yes}, said \textbf{hāṃ}. It's the letter \textbf{\hi{ह}} (h) with the \textbf{\hi{ा}} sign for long \emph{ā}, and — the key mark — the \textbf{\hi{ँ}} on top: the \textbf{chandrabindu}, the \textquotedblleft{}moon-dot.\textquotedblright{} It does one thing: it \textbf{nasalises} the vowel. So \emph{hāṃ} isn't \textquotedblleft{}hah\textquotedblright{} but a humming \emph{hãã\textsuperscript{n}}, the sound sent up through the nose, the way French does in \emph{bon}.

(Don't confuse the moon-dot \textbf{\hi{ँ}} with the plain dot \textbf{\hi{ं}} \emph{anusvara} you'll see in the next word — both nasalise, but the moon-dot rides a vowel that has nothing above it, like \emph{\hi{ा}}.)

\subsection*{You'll want to know — Saying it politely}
A bare \emph{hāṃ} is a normal, friendly yes. To an elder, a stranger, or anyone you'd show respect, Hindi adds the honorific \textbf{\hi{जी}} (\emph{jī}):

\begin{quote}
\textbf{\hi{जी} \hi{हाँ}} — \emph{jī hāṃ} — \textquotedblleft{}yes\textquotedblright{} (polite)
\end{quote}

\emph{jī} on its own is already a respectful \textquotedblleft{}yes.\textquotedblright{} You met \emph{‑jī} on names; here it warms up a plain yes into a courteous one.

\subsection*{Your turn}
Say these aloud:

\begin{itemize}
\raggedright
  \item \textquotedblleft{}hāṃ\textquotedblright{} — let the \emph{ā} hum through your nose
  \item polite — \textquotedblleft{}jī hāṃ\textquotedblright{}
  \item spot the mark — the moon-dot \textbf{\hi{ँ}} is what nasalises it
\end{itemize}

\begin{tcolorbox}[breakable,colback=teal!4,colframe=teal!35!black,title={Before you move on}]
How do you say yes in Hindi? (\textbf{hāṃ}, \hi{हाँ}.) What does the \textbf{chandrabindu} \hi{ँ} do to the vowel? (\textbf{Nasalises} it — sends it through the nose.) How do you make \emph{hāṃ} polite? (Add \textbf{\hi{जी}} — \emph{jī hāṃ}.) Next: \textbf{\hi{नहीं}}, no.
\end{tcolorbox}
\section[\hi{नहीं}]{\hi{नहीं} (nahīṃ) — no / not}

\label{lesson:HI-C07-nahin}



\begin{quote}
Hindi's \textquotedblleft{}no\textquotedblright{} carries a sound so old it also sits inside German and Latin.
\end{quote}

\subsection*{You'll want to know — The word}
\textbf{\hi{नहीं}} = \textbf{no}, said \textbf{nahīṃ} — \textbf{\hi{न}} (na) + \textbf{\hi{ह}} (h) + \textbf{\hi{ी}} (long \emph{ī}) + the \textbf{\hi{ं}} \emph{anusvara} that nasalises the end. The heart of it is that first syllable \textbf{na}, which is Hindi's ancient negator, straight from \textbf{Sanskrit na}.

\subsection*{You'll want to know — The oldest \textquotedblleft{}no\textquotedblright{} you'll meet}
That \textbf{na} is not just Sanskrit's — it is the same negative root that runs across nearly every language in this course, all descended from one ancient sound, \textbf{PIE *ne}:

\noindent
\begin{tabularx}{\linewidth}{@{}>{\raggedright\arraybackslash}X>{\raggedright\arraybackslash}X@{}}
\toprule
\textbf{from \emph{*ne}} & \textbf{} \\
\midrule
Sanskrit \emph{na} & $\to$ Hindi \textbf{nahīṃ} \\
Latin \emph{nōn} & $\to$ Spanish \textbf{no}, French \textbf{non} \\
Germanic & $\to$ German \textbf{nein}, English \textbf{no / not} \\
\bottomrule
\end{tabularx}

You met this exact idea in the German \emph{nein} lesson (\textquotedblleft{}not one\textquotedblright{}); here it is again at the other end of the family, in Hindi. The word for \emph{no} is one of the most stubbornly unchanged things human languages have.

\subsection*{You'll want to know — One word, two jobs}
Like Spanish \emph{no}, \textbf{nahīṃ} is both the \textbf{answer} and the \textbf{sentence-negator}:

\begin{itemize}
\raggedright
  \item \textbf{answer}: \emph{\hi{आप} \hi{आ} \hi{रहे} \hi{हैं}?} — \textbf{\hi{नहीं}.} (\textquotedblleft{}Are you coming?\textquotedblright{} — \textquotedblleft{}No.\textquotedblright{})
  \item \textbf{negator}: \emph{\hi{मैं} \textbf{\hi{नहीं}} \hi{जानता}} — \textquotedblleft{}I do \textbf{not} know.\textquotedblright{}
\end{itemize}

So this single word gives you both \textquotedblleft{}no\textquotedblright{} and how to make a sentence negative.

\subsection*{Your turn}
Say these aloud:

\begin{itemize}
\raggedright
  \item \textquotedblleft{}nahīṃ\textquotedblright{} — nasal at the end, from the \emph{anusvara}
  \item the ancient root — \textquotedblleft{}na… nōn… nein… no\textquotedblright{}: one PIE \emph{ne}
  \item as a negator — \textquotedblleft{}maĩ nahīṃ jāntā\textquotedblright{}, \textquotedblleft{}I don't know\textquotedblright{}
\end{itemize}

\begin{tcolorbox}[breakable,colback=teal!4,colframe=teal!35!black,title={Before you move on}]
How do you say no in Hindi? (\textbf{nahīṃ}, \hi{नहीं}.) What Sanskrit word is its core, and what ancient sound is that from? (\textbf{na}, from \textbf{PIE *ne}.) Which other course words share that root? (Latin \emph{nōn}, German \emph{nein}, English \emph{no/not} — and Spanish/French \emph{no/non}.) Besides \textquotedblleft{}no,\textquotedblright{} what else does \emph{nahīṃ} do? (It \textbf{negates a sentence} — \emph{maĩ nahīṃ jāntā}, \textquotedblleft{}I don't know.\textquotedblright{})
\end{tcolorbox}
